% GENERATED FILE. Edit canonical lessons, then run npm run generate:books.
% canonical-source-hash: fnv1a64:79dc261e5e4e4e52
% canonical-lessons: TA-C254-katan-vanku, TA-C254-celuttu, TA-C254-terntetu, TA-C254-nulai, TA-C254-veliyeru

\chapter{To borrow, To pay, To choose, To enter, To exit}
\label{ch:ta-to-borrow-to-pay-to-choose-to-enter-to-exit}
\hlchaptermodality{254}

\begin{chapteropening}
\textbf{By the end of this chapter:} \emph{I can say to borrow, to pay, to choose, to enter and to exit.}

Complete the last lesson of chapter 254: I can say to borrow, to pay, to choose, to enter and to exit.
\end{chapteropening}

\section[\texorpdfstring{\ta{கடன்} \ta{வாங்கு} (kaṭaṉ vāṅku)}{கடன் வாங்கு (kaṭaṉ vāṅku)}]{\ta{கடன்} \ta{வாங்கு} (kaṭaṉ vāṅku) — to borrow}

\label{lesson:TA-C254-katan-vanku}



\begin{quote}
Before the new one: say the Tamil for a pension, then the Tamil for a certificate.
\end{quote}

\subsection*{You'll want to know: \ta{கடன்} \ta{வாங்கு}}
\textbf{\ta{கடன்} \ta{வாங்கு}} — \emph{kaṭaṉ vāṅku} — \textquotedblleft{}to borrow\textquotedblright{}.

\begin{cousinweb}
\textbf{\ta{கடன்} \ta{வாங்கு}} (\emph{kaṭaṉ vāṅku}) is native: \emph{kaṭaṉ}, \textquotedblleft{}debt, loan\textquotedblright{} (DEDR 1113, which lists Kannada kaḍa), and \emph{vāṅku}, \textquotedblleft{}to receive, to take\textquotedblright{} (DEDR 5336): \textquotedblleft{}take a loan\textquotedblright{}.
\end{cousinweb}

\begin{grammarlens}[title={What you've built}]
One more word for this chapter.
\end{grammarlens}

\subsection*{Your turn}
\begin{itemize}
\raggedright
  \item \emph{Say it:} \emph{kaṭaṉ vāṅku}
  \item \emph{Say it:} \emph{kaṭaṉ vāṅku}, once more
  \item \emph{Say it:} say \emph{kaṭaṉ vāṅku}
  \item \emph{Recall:} say the Tamil for a pension, then the Tamil for a certificate, then say \emph{kaṭaṉ vāṅku} again
\end{itemize}

\begin{tcolorbox}[breakable,colback=teal!4,colframe=teal!35!black,title={Before you move on}]
What does \ta{கடன்} \ta{வாங்கு} mean? (\textquotedblleft{}To borrow\textquotedblright{}.) Say it once more.
\end{tcolorbox}
\section[\texorpdfstring{\ta{செலுத்து} (celuttu)}{செலுத்து (celuttu)}]{\ta{செலுத்து} (celuttu) — to pay}

\label{lesson:TA-C254-celuttu}



\begin{quote}
Before the new one: say the Tamil for a certificate, then the Tamil for to borrow.
\end{quote}

\subsection*{You'll want to know: \ta{செலுத்து}}
\textbf{\ta{செலுத்து}} — \emph{celuttu} — \textquotedblleft{}to pay\textquotedblright{}.

\begin{cousinweb}
\textbf{\ta{செலுத்து}} (\emph{celuttu}) is native Dravidian. DEDR 2781 lists Malayalam \emph{celuttuka}, \textquotedblleft{}to pay\textquotedblright{}. It is the causative of \emph{cel}, \textquotedblleft{}to go\textquotedblright{}: \textquotedblleft{}make the money go\textquotedblright{}.
\end{cousinweb}

\begin{grammarlens}[title={What you've built}]
One more word for this chapter.
\end{grammarlens}

\subsection*{Your turn}
\begin{itemize}
\raggedright
  \item \emph{Say it:} \emph{celuttu}
  \item \emph{Say it:} \emph{celuttu}, once more
  \item \emph{Say it:} say \emph{celuttu}
  \item \emph{Recall:} say the Tamil for a certificate, then the Tamil for to borrow, then say \emph{celuttu} again
\end{itemize}

\begin{tcolorbox}[breakable,colback=teal!4,colframe=teal!35!black,title={Before you move on}]
What does \ta{செலுத்து} mean? (\textquotedblleft{}To pay\textquotedblright{}.) Say it once more.
\end{tcolorbox}
\section[\texorpdfstring{\ta{தேர்ந்தெடு} (tērnteṭu)}{தேர்ந்தெடு (tērnteṭu)}]{\ta{தேர்ந்தெடு} (tērnteṭu) — to choose}

\label{lesson:TA-C254-terntetu}



\begin{quote}
Before the new one: say the Tamil for to borrow, then the Tamil for to pay.
\end{quote}

\subsection*{You'll want to know: \ta{தேர்ந்தெடு}}
\textbf{\ta{தேர்ந்தெடு}} — \emph{tērnteṭu} — \textquotedblleft{}to choose\textquotedblright{}.

\begin{cousinweb}
\textbf{\ta{தேர்ந்தெடு}} (\emph{tērnteṭu}) is native: \emph{tērntu}, \textquotedblleft{}having examined, having chosen\textquotedblright{}, from \emph{tēr} (DEDR 3419), and \emph{eṭu}, \textquotedblleft{}to take up\textquotedblright{} (DEDR 851): \textquotedblleft{}examine and take\textquotedblright{}.
\end{cousinweb}

\begin{grammarlens}[title={What you've built}]
One more word for this chapter.
\end{grammarlens}

\subsection*{Your turn}
\begin{itemize}
\raggedright
  \item \emph{Say it:} \emph{tērnteṭu}
  \item \emph{Say it:} \emph{tērnteṭu}, once more
  \item \emph{Say it:} say \emph{tērnteṭu}
  \item \emph{Recall:} say the Tamil for to borrow, then the Tamil for to pay, then say \emph{tērnteṭu} again
\end{itemize}

\begin{tcolorbox}[breakable,colback=teal!4,colframe=teal!35!black,title={Before you move on}]
What does \ta{தேர்ந்தெடு} mean? (\textquotedblleft{}To choose\textquotedblright{}.) Say it once more.
\end{tcolorbox}
\section[\texorpdfstring{\ta{நுழை} (nuḻai)}{நுழை (nuḻai)}]{\ta{நுழை} (nuḻai) — to enter}

\label{lesson:TA-C254-nulai}



\begin{quote}
Before the new one: say the Tamil for to pay, then the Tamil for to choose.
\end{quote}

\subsection*{You'll want to know: \ta{நுழை}}
\textbf{\ta{நுழை}} — \emph{nuḻai} — \textquotedblleft{}to enter\textquotedblright{}.

\begin{cousinweb}
\textbf{\ta{நுழை}} (\emph{nuḻai}) is native Dravidian. DEDR 3714 lists Malayalam \emph{nuḻayuka}, \textquotedblleft{}to creep through\textquotedblright{}.
\end{cousinweb}

\begin{grammarlens}[title={What you've built}]
One more word for this chapter.
\end{grammarlens}

\subsection*{Your turn}
\begin{itemize}
\raggedright
  \item \emph{Say it:} \emph{nuḻai}
  \item \emph{Say it:} \emph{nuḻai}, once more
  \item \emph{Say it:} say \emph{nuḻai}
  \item \emph{Recall:} say the Tamil for to pay, then the Tamil for to choose, then say \emph{nuḻai} again
\end{itemize}

\begin{tcolorbox}[breakable,colback=teal!4,colframe=teal!35!black,title={Before you move on}]
What does \ta{நுழை} mean? (\textquotedblleft{}To enter\textquotedblright{}.) Say it once more.
\end{tcolorbox}
\section[\texorpdfstring{\ta{வெளியேறு} (veḷiyēṟu)}{வெளியேறு (veḷiyēṟu)}]{\ta{வெளியேறு} (veḷiyēṟu) — to exit, to go out}

\label{lesson:TA-C254-veliyeru}



\begin{quote}
Before the new one: say the Tamil for to choose, then the Tamil for to enter.
\end{quote}

\subsection*{You'll want to know: \ta{வெளியேறு}}
\textbf{\ta{வெளியேறு}} — \emph{veḷiyēṟu} — \textquotedblleft{}to exit, to go out\textquotedblright{}.

\begin{cousinweb}
\textbf{\ta{வெளியேறு}} (\emph{veḷiyēṟu}) is native: \emph{veḷi}, \textquotedblleft{}outside\textquotedblright{} (DEDR 5498), and \emph{ēṟu}, \textquotedblleft{}to rise, to go up\textquotedblright{} (DEDR 916): \textquotedblleft{}rise to the outside\textquotedblright{}.
\end{cousinweb}

\begin{grammarlens}[title={What you've built}]
One more word for this chapter.
\end{grammarlens}

\subsection*{Your turn}
\begin{itemize}
\raggedright
  \item \emph{Say it:} \emph{veḷiyēṟu}
  \item \emph{Say it:} \emph{veḷiyēṟu}, once more
  \item \emph{Say it:} say \emph{veḷiyēṟu}
  \item \emph{Recall:} say the Tamil for to choose, then the Tamil for to enter, then say \emph{veḷiyēṟu} again
\end{itemize}

\begin{tcolorbox}[breakable,colback=teal!4,colframe=teal!35!black,title={Before you move on}]
What does \ta{வெளியேறு} mean? (\textquotedblleft{}To exit, to go out\textquotedblright{}.) Say it once more.
\end{tcolorbox}
